\slajdid{10-s091}
\begin{frame}[label=10-s091]
\gusttekst\setlength{\abovedisplayskip}{3pt}\setlength{\belowdisplayskip}{3pt}\setlength{\abovedisplayshortskip}{2pt}\setlength{\belowdisplayshortskip}{2pt}\begin{columns}[T,onlytextwidth]\column{.35\textwidth}\centering\input{slike/tikz/s090_ecl_odvojena_napajanja.tex}\column{.63\textwidth}
Ulazni strujni kapaciteti su
\[I_{IL}=\frac{V_I-V_{EE}}{R_B}=\frac{V_{IL}-V_{EE}}{R_B}\]
\[I_{IH}=\frac{V_I-V_{EE}}{R_B}+\frac{I_E}{\beta_F}=\frac{V_{OH}-V_{EE}}{R_B}+\frac{I_E}{\beta_F}\]
Dok su strujni kapaciteti izlaza
\[I_{OH}=\beta_F\frac{V_{CC}-(V_{OHmin}+V_{BE})}{R_C}\]
pri čemu je $V_{OHmin}$ određeno time da tranzistor T1 (ili T2) ne odu u zasićenje
\par\bigskip Strujni kapacitet logičke nule je praktično beskonačan, pošto je tada odgovarajući tranzistor zakočen pa struja baze baferskog tranzistora ne utiče na njegov režim rada.
\end{columns}
\end{frame}
